\chapter*{General Conclusion}
\addcontentsline{toc}{chapter}{General Conclusion}

This end-of-studies project was driven by a real operational need at KTS: replacing a paper-based production tracking process with an integrated digital platform. Over the course of this report, we progressed from problem analysis to technical design and then to the practical implementation of a web application designed to better structure daily production activities and improve coordination among stakeholders.

The developed solution, \textbf{Production Tracker}, addresses the core business workflows identified during the analysis phase. It provides role-based authentication, work order management across the main production stages, stock monitoring and reservation, real-time notifications, and instant communication through chat. By bringing these functions together in a single system, the platform reduces information dispersion, limits manual errors, and improves the traceability of production operations.

From a technical perspective, the chosen stack (Next.js, Node.js/Express, PostgreSQL, Prisma, and Socket.IO) proved well aligned with the project objectives. It made it possible to build a modular architecture, ensure a clear separation of responsibilities between frontend and backend, and develop both transactional and real-time features efficiently. The resulting system offers a solid foundation for maintainability and future evolution.

Beyond the technical implementation, this work brings clear organizational value to the host company. Administrators now have better visibility over production progress, workers receive clearer task guidance, and communication delays are reduced through shared real-time information. In this way, the project contributes directly to KTS's digital transformation effort and supports more reliable day-to-day decision making.

Like any first complete version, the current platform can still be improved. Future work may include advanced analytics dashboards, richer KPI reporting, tighter integration with ERP systems, mobile-first access for shop-floor usage, and additional automation for planning and forecasting. These extensions would strengthen strategic monitoring while preserving the operational gains already achieved.

In conclusion, this internship project turned a real industrial problem into a functional and scalable software solution. It also marked an important personal and professional milestone, bringing together software engineering practices, system design, and iterative delivery in a real business context.
